\chapter{The Representations of \texorpdfstring{$\GL(2)$}{GL(2)} over Local and Global Fields}

\label{ch:4}
\label{ch:4}

In this section we shall describe some basic features of the theory of representations for $\GL(2)$ over local and global fields.

\section{The Archimedean Places}\label{sec:4a}

We deal first with $\GL(2,\R)$ since our primary interest is in automorphic forms defined over $\Q$. The theory for $\GL(2,\C)$ is simpler, requires little additional attention, and is included for the sake of completeness.

In Section~\ref{ch:2} we have described the irreducible unitary representations of $\SL(2,\R)$. Our present point of departure is the well-known fact that any such representation $\pi$ is ``admissible'' in the following sense.

Consider the universal enveloping algebra of the complexification of the Lie algebra of $\SL(2,\R)$. If $\pi$ is realized on the Hilbert space $H$, and $\sigma$ is an irreducible representation of $K = \SO(2)$, let $H(\sigma)$ denote the subspace of vectors $v$ in $H$ which transform according to $\sigma$. (Thus $\pi(K)v$ spans a finite-dimensional subspace of $H$ equivalent to a finite number of copies of $\sigma$.) The algebraic direct sum

\[
H^0 = \bigoplus_{\sigma} H(\sigma)
\]

is dense in $H$ and comprises the subspace of $K$-finite vectors. The crucial fact is that by differentiation $\pi$ induces a representation of the enveloping algebra on the space $H^0$. This representation is (algebraically) irreducible since the unitary representation $\pi$ is (topologically) irreducible. It is ``admissible'' in the sense that its restriction to the Lie algebra of $K$ decomposes into finite-dimensional representations with finite multiplicities. (Equivalently, the dimension of $H(\sigma)$ is finite for each $\sigma$.)

Following Harish-Chandra and Jacquet-Langlands we henceforth focus our attention not on irreducible unitary representations of $G = \GL(2,\R)$ but rather on irreducible admissible representations of an appropriate group algebra of $G$.

\begin{definition}\label{def:4.1}

Let $\mathfrak{G}$ denote the universal enveloping algebra of the complexification of the Lie algebra of $G$. Let $\varepsilon_{-}$ denote the Dirac measure at $\begin{pmatrix} -1 & 0 \\ 0 & 1 \end{pmatrix}$. Then the Hecke group algebra of $G$ is by definition the direct sum

\[
\mathcal{H}(G) = \mathfrak{G} \oplus (\varepsilon_{-}) * \mathfrak{G}.
\]

\end{definition}

Observe that $\mathcal{H}(G)$ is an algebra (under convolution product) of distributions supported in the subgroup $\left\{\begin{pmatrix} \pm 1 & 0 \\ 0 & 1 \end{pmatrix}\right\}$ of $G$. It is the simplest substitute for the enveloping algebra of $G$ which takes into account the fact that $G$ has two connected components.

\begin{definition}\label{def:4.2}

Suppose $\pi$ is a representation of the algebra $\mathcal{H}(G)$ on a complex vector space $V$. Then $\pi$ is said to be admissible if its restriction to the Lie algebra of $K = O(2,\R)$ decomposes into finite-dimensional representations with finite multiplicities.

\end{definition}

\begin{remark}\label{rem:4.3}

Almost all the irreducible representations of $K = O(2,\R)$ are two dimensional. Suppose $\sigma$ is any such representation of (the Lie algebra of) $K$ and $V(\sigma)$ denotes the subspace of vectors in $V$ which transform according to some multiple of $\sigma$. Then admissibility of $(\pi, V)$ amounts to the assertion that

\begin{equation}\label{eq:4.1}
V = \bigoplus_{\sigma} V(\sigma) \quad \text{(algebraic sum)}\tag{4.1}
\end{equation}

with each $V(\sigma)$ finite dimensional.

\end{remark}

Our task now is to describe all the irreducible admissible representations up to equivalence. This task is simplified by the fact that all such representations are known to be subquotients of the following basic representations of $\mathcal{H}(G)$, the so-called principal series representations.

Let $\mu_1, \mu_2$ denote any quasi (i.e.\ not necessarily unitary) characters of $\R^{\times}$. Let $\mathcal{H}(\mu_1, \mu_2)$ denote the vector space of functions $\phi(g)$ on $G$ which are right $K$-finite (cf.\ the definition in Section~\ref{ch:2}A) and such that

\begin{equation}\label{eq:4.2}
\phi\!\left(\begin{pmatrix} t_1 & * \\ 0 & t_2 \end{pmatrix} g\right) = \mu_1(t_1)\mu_2(t_2)\left|\frac{t_1}{t_2}\right|^{1/2}\phi(g)\tag{4.2}
\end{equation}

for all $t_1, t_2 \in \R^{\times}$. According to the Iwasawa decomposition,

\begin{equation}\label{eq:4.3}
G = NA\,\SO(2)\tag{4.3}
\end{equation}

(where $A = \left\{\begin{pmatrix} t_1 & 0 \\ 0 & t_2 \end{pmatrix}\right\}$). Therefore, each $\phi$ in $\mathcal{H}(\mu_1, \mu_2)$ is completely determined by its restriction to $\SO(2,\R)$ (and in particular is infinitely differentiable since this restriction must be a trigonometric polynomial).

If $X$ belongs to the Lie algebra we define (as usual)

\begin{equation}\label{eq:4.4}
\phi * X(g) = \frac{\dd}{\dd t} \phi(g \exp(-tX))\big|_{t=0}\tag{4.4}
\end{equation}

and let $\check{X}$ denote $-X$. By $\rho(\mu_1, \mu_2)$ we denote the representation of $\mathcal{H}(G)$ on $\mathcal{H}(\mu_1, \mu_2)$ determined by

\begin{equation}\label{eq:4.5}
\rho(\mu_1, \mu_2)(X)\phi = \phi * \check{X}.\tag{4.5}
\end{equation}

This representation is the lift to the enveloping algebra of the representation of $G$ induced from the character

\[
\begin{pmatrix} t_1 & * \\ 0 & t_2 \end{pmatrix} \to \mu_1(t)\mu_2(t)
\]

of $B$. It is admissible precisely because each $\phi$ in $\mathcal{H}(\mu_1, \mu_2)$ restricts to a trigonometric polynomial on $\SO(2,\R)$. What is remarkable is that every irreducible admissible representation of $\mathcal{H}(G)$ is a subquotient of some such $\rho(\mu_1, \mu_2)$. Therefore the classification of irreducible admissible representations is reduced to the study of how (and when) these $\rho(\mu_1, \mu_2)$ decompose.

To analyze the reducibility of $\rho(\mu_1, \mu_2)$ one can compute the action of certain Lie algebra elements on convenient basis elements of $\mathcal{H}(\mu_1, \mu_2)$. To this end, write $\mu_i(t) = |t|^{s_i}[\operatorname{sgn}(t)]^{m_i}$, where $m_i = 0$ or $1$, and $s_i \in \C$. Then

\begin{equation}\label{eq:4.6}
\mu(t) = \mu_1\mu_2^{-1}(t) = |t|^s[\operatorname{sgn}(t)]^m\tag{4.6}
\end{equation}

with $s = s_1 - s_2$ and $m = m_1 - m_2$. For each $n \equiv m \pmod{2}$ the function

\begin{equation}\label{eq:4.7}
\phi_n\!\left(\begin{pmatrix} t_1 & * \\ 0 & t_2 \end{pmatrix} \begin{pmatrix} \cos\theta & -\sin\theta \\ \sin\theta & \cos\theta \end{pmatrix}\right) = \mu_1(t_1)\mu_2(t_2)\left|\frac{t_1}{t_2}\right|^{1/2} e^{-in\theta}\tag{4.7}
\end{equation}

belongs to $\mathcal{H}(\mu_1, \mu_2)$ and the set $\{\phi_n\}$ is obviously a basis for $\mathcal{H}(\mu_1, \mu_2)$. The relevant elements of $\mathfrak{G}$ are

\[
U = \begin{pmatrix} 0 & 1 \\ -1 & 0 \end{pmatrix} \quad J = \begin{pmatrix} 1 & 0 \\ 0 & 1 \end{pmatrix} \quad V_+ = \begin{pmatrix} 1 & i \\ i & -1 \end{pmatrix} \quad V_- = \begin{pmatrix} 1 & -i \\ -i & -1 \end{pmatrix} \quad Z = \begin{pmatrix} 1 & 0 \\ 0 & -1 \end{pmatrix}
\]

\[
X_+ = \begin{pmatrix} 0 & 1 \\ 0 & 0 \end{pmatrix} \quad X_- = \begin{pmatrix} 0 & 0 \\ 1 & 0 \end{pmatrix}
\]

as well as the Casimir operator

\[
D = X_+X_- + X_-X_+ + \frac{Z^2}{2}
\]

which belongs to (the center of) $\mathfrak{G}$. Writing $\rho$ for $\rho(\mu_1, \mu_2)$, we have

\begin{enumerate}[(i)]

\item $\rho(U)\phi_n = in\phi_n$

\item $\rho(V_+)\phi_n = (s + 1 + n)\phi_{n+2}$

\item $\rho(V_-)\phi_n = (s + 1 - n)\phi_{n-2}$

\item $\rho(\varepsilon_{-})\phi_n = (-1)^{m_1}\phi_{-n}$

\item $\rho(D)\phi_n = \frac{s^2 - 1}{2}\phi_n$, and

\item $\rho(J)\phi_n = (s_1 + s_2)\phi_n$.

\end{enumerate}

From this it follows that

\[
\rho(V_+)^p\phi_n = (s + 1 + n) \cdots (s + 2p - 1 + n)\phi_{n+2p}
\]

and

\[
\rho(V_-)^p\phi_n = (s + 1 - n) \cdots (s + 2p - 1 - n)\phi_{n-2p}.
\]

Now the elements $U, Z, V_+$ and $V_-$ form a $\C$-basis for the Lie algebra of $G$ and the eigenvectors of $\rho(U)$ span $\mathcal{H}(\mu_1, \mu_2)$. Consequently the above computations show that $\rho(\mu_1, \mu_2)$ is reducible if and only if $s - m$ is an odd integer, i.e.\ if and only if $\mu(t) = t^p \operatorname{sgn}(t)$ for some integer $p$. More precisely, we have:

\begin{theorem}\label{thm:4.4}

\begin{enumerate}[(a)]

\item If $\mu_1\mu_2^{-1}$ is not of the form $t \to t^p\operatorname{sgn}(t)$, with $p$ a non-zero integer, then $\mathcal{H}(\mu_1, \mu_2)$ is irreducible under the action of $\mathcal{H}(G)$;

\item If $\mu_1\mu_2^{-1}(t) = t^p\operatorname{sgn}(t)$, with $p > 0$, then $\mathcal{H}(\mu_1, \mu_2)$ contains exactly one invariant subspace, namely,

\[
\mathcal{H}^s(\mu_1, \mu_2) = \{\dots, \phi_{-p-3}, \phi_{-p-1}, \phi_{p+1}, \phi_{p+3}, \dots\}
\]

and the quotient $\mathcal{H}^f(\mu_1, \mu_2) = \mathcal{H}(\mu_1, \mu_2)/\mathcal{H}^s(\mu_1, \mu_2)$ is finite dimensional;

\item If $\mu(t) = t^p\operatorname{sgn}(t)$ with $p < 0$, then the only invariant subspace of $\mathcal{H}(\mu_1, \mu_2)$ is

\[
\mathcal{H}^f(\mu_1, \mu_2) = \{\phi_{p+1}, \phi_{p+3}, \dots, \phi_{-p-3}, \phi_{-p-1}\}.
\]

\end{enumerate}

\end{theorem}

The representation $\rho(\mu_1, \mu_2)$ will be denoted by $\pi(\mu_1, \mu_2)$ if it is irreducible; in case it is not, the obvious representation on the finite dimensional space $\mathcal{H}^f(\mu_1, \mu_2)$ will still be denoted by $\pi(\mu_1, \mu_2)$. The representation on the infinite dimensional subspace (or quotient) $\mathcal{H}^s(\mu_1, \mu_2)$ will be denoted by $\sigma(\mu_1, \mu_2)$ and viewed as a member of the discrete series for $G$. It is defined only when $\mu(t) = t^p\operatorname{sgn}(t)$ for some non-zero integer $p$. The representations $\pi(\mu_1, \mu_2)$ exhaust the so-called principal series for $G$.

\begin{theorem}\label{thm:4.5}

Every irreducible admissible representation of $\mathcal{H}(G)$ is either a $\pi(\mu_1, \mu_2)$ or a $\sigma(\mu_1, \mu_2)$; the only equivalences between these representations are the following:

\[
\pi(\mu_1, \mu_2) \approx \pi(\mu_2, \mu_1)
\]

and

\[
\sigma(\mu_1, \mu_2) \approx \sigma(\mu_2, \mu_1) \approx \sigma(\mu_1\eta, \mu_2\eta) \approx \sigma(\mu_2\eta, \mu_1\eta),
\]

where $\eta(t) = \operatorname{sgn}(t)$.

\end{theorem}

\begin{remark}\label{rem:4.6}

The irreducible admissible representations of $\mathcal{H}(G)$ are often indexed by two particularly convenient parameters. For the (infinite dimensional) principal series representation $\pi(\mu_1, \mu_2)$ these parameters are $s = s_1 - s_2$ and $t = s_1 + s_2$ (cf.\ 4.5). Note then that two principal series representations $\pi$ and $\pi'$ will share the same eigenvalue for the Casimir operator if and only if $|s| = |s'|$. (They will share the same central character if and only if $t = t'$.) If $\pi$ is equivalent to the discrete series representation $\sigma(\mu_1, \mu_2)$ then the convenient parameters for $\pi$ are (again) $p = s_1 - s_2$ and $t = s_1 + s_2$ (if $\mu_i(t) = |t|^{s_i}[\operatorname{sgn}(t)]^{m_i}$). The significance of $p$ is that the space of $\pi(p, t)$ will contain the functions $\phi_{p+1}, \phi_{p+3}, \cdots$, but not the function $\phi_{p-1}$. Consequently $\phi_{p+1}$ may be viewed as the lowest weight vector for $\pi(p, t)$.

\end{remark}

\begin{remark}\label{rem:4.7}

It can be shown that a given principal series representation $\pi(\mu_1, \mu_2)$ of $\mathcal{H}(G)$ corresponds to a unitary representation of $G$ if and only if either both $\mu_1$ and $\mu_2$ are unitary, i.e.\ the parameters $s$ and $t$ just introduced are both pure imaginary, or just $t$ is imaginary and $s = s_1 - s_2$ is real, non-zero, and between $-1$ and $1$. In accordance with the terminology already introduced for $\SL(2,\R)$, $\pi(\mu_1, \mu_2)$ is called a continuous series representation when this first possibility obtains, and a complementary series representation when the second does.

The discrete series representations $\sigma(\mu_1, \mu_2)$ are similarly ``unitarizable'' if and only if their central character is unitary, i.e.\ if and only if $t = s_1 + s_2$ is pure imaginary (cf.\ 4.7(vi)).

\end{remark}

\begin{remark}\label{rem:4.8}

(\ensuremath{\GL(2,\C)}.) Since $\GL(2,\C)$ is connected it is reasonable to expect that in this case the correspondence between ``nice'' representations of $G$ and ``nice'' representations of the Lie algebra of $G$ will be one-to-one. Therefore for $\GL(2,\C)$ the group algebra $\mathcal{H}(G)$ is defined to be precisely the enveloping algebra of the (complexification of the) Lie algebra of $G$. The ``principal series'' $\rho(\mu_1, \mu_2)$ are introduced in the obvious way for each pair of quasi-characters $\mu_1, \mu_2$ of $\C^{\times}$ and again it is the case that every irreducible admissible representation of $\mathcal{H}(G)$ is a subquotient of some $\rho(\mu_1, \mu_2)$. Furthermore, $\rho(\mu_1, \mu_2)$ is irreducible except when $\mu_1\mu_2^{-1}(z) = z^p\overline{z}^q$, with $p, q \in \Z$ and $pq > 0$, in which case $\mathcal{H}(\mu_1, \mu_2)$ ``contains'' exactly one infinite dimensional irreducible subquotient. However, in contrast to the real case, this irreducible subquotient is actually equivalent to some irreducible $\rho(\mu_1', \mu_2')$ for appropriate $\mu_1', \mu_2'$! In particular, there are no discrete series for $\GL(2,\C)$.

\end{remark}

\section{The \texorpdfstring{$p$}{p}-adic Theory}\label{sec:4b}

Our purpose is to describe the classification of all irreducible admissible representations of

\[
G = \GL(2,F)
\]

when $F$ is a non-archimedean locally compact field. We also shall collect some basic facts concerning these representations. These facts and others play a crucial role in the theory of automorphic forms.

\subsection{Admissibility}

We start with some notation. Throughout this section, $F$ will be a finite extension of the $p$-adic number field or a field of formal power series in one variable over a finite field. The symbol $O_F$ will denote the ring of integers of $F$. The absolute value on $F$ is defined by the relation $\dd(ax) = |a|\dd x$ where $\dd x$ is an invariant measure on the additive group of $F$. The prime ideal of $O_F$, defined by $|a| < 1$, will be denoted by $P_F$. It is generated (say) by $\varpi$, and $O_F/P_F$ is a finite field with $q$ elements, $q$ some power of $p$. The standard maximal compact (open) subgroup of $G$ is

\[
K = \GL(2, O_F).
\]

Now $G$ has no enveloping algebra of differential operators but there is a natural group algebra for $G$ which is fundamental for its representation theory, namely the (Hecke) group algebra $\mathcal{H}(G)$ consisting of all locally constant compactly supported functions on $G$. (This is an algebra for the group convolution product

\[
f * g(x) = \int_G f(xy^{-1})g(y)\dd^{*}y,
\]

where $\dd^{*}y$ denotes the Haar measure for $G$ which assigns the measure 1 to $K$.)

Following Jacquet and Langlands, we say that a representation $\pi$ of $\mathcal{H}(G)$ on a complex vector space $V$ is admissible if for every $v$ in $V$ there is an $f$ in $\mathcal{H}(G)$ such that $\pi(f)v = v$, and if every $\pi(f)$ maps $V$ onto a finite dimensional space. This definition of admissibility is motivated by the following considerations.

Suppose $\pi$ is an irreducible unitary representation of $G$ on some Hilbert space $H$. Then $\pi$ induces a representation of $\mathcal{H}(G)$ through the formula

\[
\pi(f) = \int_G f(g)\pi(g)\dd^{*}g,
\]

and, as in the real case, $\pi(f)$ defines a representation of $\mathcal{H}(G)$ in the subspace of $K$-finite functions of $H$ which is admissible in the above sense. However, in contrast to the real case, $\pi(g)$ itself acts in the space of $K$-finite vectors. Therefore it should not seem surprising that for $p$-adic groups (and admissible representations) the correspondence between representations of $G$ and representations of the group algebra of $G$ is completely transparent. To see this, we make the following definition.

\begin{definition}\label{def:4.9}

Suppose $\pi$ is a representation of $G$ on a complex vector space $V$ (the space of $K$-finite vectors of some unitary $G$-space, for example). Then $\pi$ is said to be admissible if (i) the stabilizer in $G$ of each $v$ in $V$ is an open subgroup of $K$, and (ii) the subspace of $V$ fixed by any open subgroup of $K$ is finite-dimensional.

\end{definition}

\begin{remark}\label{rem:4.10}

In Definition~\ref{def:4.9} no topology is imposed on the space $V$ and so Condition (i) is meant as a substitute for the usual continuity condition for $\pi$. Condition (ii) is really the crucial assumption of admissibility. It implies that each operator

\[
\pi(f) = \int_G f(g)\pi(g)\dd^{*}g
\]

is of finite rank and therefore that the character of $\pi$, defined by

\[
\theta_{\pi}(f) = \operatorname{Trace}\pi(f),
\]

always exists (at least as a ``distribution'' on $G$).

\end{remark}

\begin{lemma}\label{lem:4.11}

If $\pi(g)$ is an admissible representation of $G$, then $\pi(f)$ defines an admissible representation of $\mathcal{H}(G)$. More significantly, if $\pi(f)$ is an admissible representation of $\mathcal{H}(G)$, then there exists an admissible representation $\pi$ of $G$ such that

\[
\pi(f) = \int_G f(g)\pi(g)\dd^{*}g.
\]

\end{lemma}

The proof of this lemma is tedious but straightforward. Since it is to be found in Section~\ref{ch:2} of \cite{JacquetLanglands} we simply content ourselves with the observation that for each $v$ in $V$, the integral

\[
\pi(f)v = \int_G f(g)\pi(g)v\dd^{*}g
\]

is actually a finite sum. Indeed, since $f$ belongs to $\mathcal{H}(G)$, Condition (i) of Definition~\ref{def:4.9} implies that $f(g)\pi(g)v$ takes on only finitely many values.

\begin{remark}\label{rem:4.12}

The correspondence (established in Lemma~\ref{lem:4.11}) between admissible representations of $G$ and $\mathcal{H}(G)$ is more than simply one-to-one. For if $V_1$ is a subspace of $V$ invariant for $G$ it is also invariant for $\mathcal{H}(G)$, and vice versa. Similarly, an operator in $V$ commutes with the action of $G$ if and only if it commutes with the action of $\mathcal{H}(G)$.

\end{remark}

Henceforth we shall deal exclusively with irreducible admissible representations of $\mathcal{H}(G)$. Before describing the classification of such representations we offer yet another useful definition of admissibility.

\begin{definition}\label{def:4.13}

Suppose $\pi$ is a representation of $G$ on a complex vector space $V$. Then the restriction of $\pi$ to $K$ contains any given irreducible representation of $K$ at most finitely many times. Equivalently, if $V(\sigma)$ denotes the subspace of $v$ in $V$ which transform according to the continuous irreducible representation $\sigma$ of $K$, then

\begin{enumerate}[(i)]

\item $V = \bigoplus_{\sigma} V(\sigma)$ (algebraic)

\item the dimension of each $V(\sigma)$ is finite.

\end{enumerate}

\end{definition}

This definition is entirely analogous to the usual definition for Lie groups. Its equivalence with Definition~\ref{def:4.9} seems never to have been verified in print and therefore we write out the details of this argument below:

\begin{lemma}\label{lem:4.14}

Definitions \ref{def:4.9} and \ref{def:4.13} are equivalent.

\end{lemma}

\begin{proof}

Suppose first that $\pi$ is admissible in the sense of Definition~\ref{def:4.9}. To prove that

\begin{equation}\label{eq:4.8}
V = \bigoplus_{\sigma} V(\sigma)\tag{4.8}
\end{equation}

we consider the stabilizer $U$ in $K$ of any vector $v$ in $V$. By Condition (i) of Definition~\ref{def:4.9} this $U$ is open, and therefore we may find $k_1, \cdots, k_N$ such that $K = \bigcup_{i=1}^{N} k_iU$. Since this implies that the space of translates of $v$ under $K$ is finite-dimensional, \eqref{eq:4.8} follows by the arbitrariness of $v$. To prove that each $V(\sigma)$ is finite-dimensional we recall that each continuous irreducible representation of $K$ is finite-dimensional. Consequently, there is an open subgroup $K^{*}$ of $K$ such that $\sigma(k)$ is the identity operator for each $k \in K^{*}$. (Indeed, if $\{v_1, \cdots, v_d\}$ is any basis for the space of $\sigma$ in $V$, $K^{*}$ may be taken to be the intersection of the stability groups of these $v_i$.) But this implies that $V(\sigma)$ is contained in the space of fixed vectors for $K^{*}$. Therefore $V(\sigma)$ is finite-dimensional by Condition (ii) of Definition~\ref{def:4.9}.

Now suppose that $\pi$ is admissible in the sense of Definition~\ref{def:4.13} and fix any $v$ in $V$. By Condition (i) and (ii) of Definition~\ref{def:4.13} we may assume (without any loss of generality) that $v$ belongs to an irreducible subspace of some finite-dimensional space $V(\sigma)$. But then the stability subgroup of $v$ contains some open subgroup of $K$ by the continuity of $\sigma$. Therefore it remains only to verify that Condition (ii) of Definition~\ref{def:4.9} holds.

To this end let $U$ denote an arbitrary open subgroup of $K$ and $V_U$ the corresponding subspace of $V$ fixed by $U$. Without loss of generality we may assume that $U$ is normal in $K$. To prove that $V_U$ is finite-dimensional it will suffice to prove that $V_U$ is contained in the direct sum of a finite number of spaces $V(\sigma_i)$ (since each of these spaces is finite-dimensional by Condition (ii) of Definition~\ref{def:4.13}). But now we claim that only finitely many representations $\sigma$ contain the identity representation upon restriction to $U$ and that the direct sum of the spaces in $V$ fixed by these representations contains $V_U$. Indeed if $\sigma_i$ is such a representation, and $w$ is any $U$-invariant vector in the space of $\sigma_i$, then $\pi(K)w$ is also $U$-invariant (since $\pi(u)\pi(k)w = \pi(k)w$ if and only if $\pi(k^{-1})\pi(u)\pi(k)w = w$, an obvious identity by the normality of $U$ in $K$). Therefore there is a one-to-one correspondence between such $\sigma$ and equivalence class of irreducible representations of the finite group $K/U$ (on $V(\sigma) \cap V_U$). From this our claim (and hence the lemma) follows.

\end{proof}

\begin{remark}\label{rem:4.15}

(Schur's Lemma.) Schur's Lemma holds for all irreducible admissible representations, unitary or not. Indeed, any linear operator $A$ in $V$ which commutes with $\pi$ also operates in every finite dimensional subspace $\pi(f)(V)$ of $V$. Therefore $A$ must have at least one eigenvector $v_0$. Since $\pi(g)v_0$ is again an eigenvector for $A$ (for any $g \in G$) it follows from the irreducibility of $\pi$ that $A$ is a scalar.

One consequence of Schur's Lemma is that any irreducible admissible representation of $G$ defines a character of the center $Z$ of $G$ (the so-called central character). Another is that every finite-dimensional irreducible admissible representation of $\GL(2)$ (over a $p$-adic field) is one dimensional and of the form $\chi(\det g)$ for some quasi-character $\chi$ of $F^{\times}$. Indeed, if the space of $\pi$ is finite-dimensional, its kernel is an open (hence non-trivial) normal subgroup of $G$. Now this kernel certainly contains

\[
\begin{pmatrix} 1 & x \\ 0 & 1 \end{pmatrix} = \begin{pmatrix} a^{-1} & 0 \\ 0 & 1 \end{pmatrix} \begin{pmatrix} 1 & ax \\ 0 & 1 \end{pmatrix} \begin{pmatrix} a & 0 \\ 0 & 1 \end{pmatrix}
\]

and $\begin{pmatrix} 1 & 0 \\ x & 1 \end{pmatrix}$ since $|ax|$ can be made arbitrarily small for $a$ sufficiently small. Therefore the kernel of $\pi$ contains the group generated by the elements $\begin{pmatrix} 1 & x \\ 0 & 1 \end{pmatrix}$ and $\begin{pmatrix} 1 & 0 \\ x & 1 \end{pmatrix}$, namely $\SL(2,F)$! So since $G = \left\{\begin{pmatrix} \det g & 0 \\ 0 & 1 \end{pmatrix}\right\}\SL(2,F)$ our claim follows from Schur's Lemma.

\end{remark}

\subsection{Classification of Admissible Representations}

The main classification result is that every irreducible admissible representation $\pi$ of $G$ is a subrepresentation of some principal series representation unless $\pi$ is a so-called supercuspidal representation characterized by the fact that its matrix coefficients are compactly supported modulo the center. This is to be contrasted with the real case where no such supercuspidal representations exist.

To describe the principal series, let $\mu_1, \mu_2$ be any two quasi-characters of $F^{\times}$ and let $\mathcal{H}(\mu_1, \mu_2)$ denote the space of all locally constant functions $\phi$ on $G$ such that

\begin{equation}\label{eq:4.9}
\phi\!\left(\begin{pmatrix} t_1 & * \\ 0 & t_2 \end{pmatrix} g\right) = \mu_1(t_1)\mu_2(t_2)|t_1/t_2|^{1/2}\phi(g)\tag{4.9}
\end{equation}

for all $t_1, t_2 \in F^{\times}$. The group $G$ acts on $\mathcal{H}(\mu_1, \mu_2)$ through right translations and the resulting representation of $G$ is called a principal series representation (at least when it is irreducible) and denoted by $\rho(\mu_1, \mu_2)$. Each such representation is clearly admissible. Indeed the Iwasawa decomposition

\begin{equation}\label{eq:4.10}
G = BK\tag{4.10}
\end{equation}

implies that functions in $\mathcal{H}(\mu_1, \mu_2)$ are completely determined by their restriction to $K$. In fact $\mathcal{H}(\mu_1, \mu_2)$ is naturally isomorphic to the space of locally constant functions on $K$ such that

\[
\phi\!\left(\begin{pmatrix} u_1 & x \\ 0 & u_2 \end{pmatrix} k\right) = \mu_1(u_1)\mu_2(u_2)\phi(k)
\]

for all $u_1, u_2 \in O_F^{\times}$ and $x \in O_F$. Therefore, if $U$ is any open subgroup of $K$, the restriction of a function invariant under $U$ is a function on the finite set $K/U$. Thus the space of all such functions is finite-dimensional (cf.\ Definition~\ref{def:4.9}).

It is important to point out that $\rho(\mu_1, \mu_2)$ is nothing more or less than the representation of $G$ induced from the one-dimensional representation

\[
\begin{pmatrix} t_1 & x \\ 0 & t_2 \end{pmatrix} \to \mu_1(t_1)\mu_2(t_2)
\]

of $B = \left\{\begin{pmatrix} t_1 & x \\ 0 & t_2 \end{pmatrix}: t_i \in F^{\times}, x \in F\right\}$.

\begin{definition}\label{def:4.16}

Suppose $\pi$ is any admissible representation of $G$ on $V$. Let $\tilde{V}$ denote the subspace of those $v^{*}$ in $V^{*}$ (the algebraic dual of $V$) which are invariant under some open subgroup of $G$. Then the contragredient of $\pi$, denoted by $\tilde{\pi}(g)$, is the restriction of ${}^{t}\pi(g^{-1})$ to $\tilde{V}$.

\end{definition}

Because

\[
\tilde{V} = \bigoplus_{\sigma} V(\sigma)^{*}
\]

it is clear that $\tilde{\pi}$ is admissible as soon as $\pi$ is. Furthermore, if $V_1$ is an invariant subspace of $V$ for $\pi$, then its ``orthogonal complement''

\[
(V_1)^{\perp} = \{v^{*} \in V: v^{*}(v) = 0 \text{ for all } v \in V_1\}
\]

is invariant for $\tilde{\pi}$ and $(V_1^{\perp})^{\perp} = V_1$. Therefore there is a one-to-one correspondence between invariant subspaces of $V$ and $\tilde{V}$. Note that $(\tilde{\pi}(g)\tilde{v}, \pi(g)v) = (\tilde{v}, v)$ if $(\cdot, \cdot)$ denotes the duality between $V$ and $\tilde{V}$. Note also that $(\tilde{\pi})^{\sim} = \pi$. This notion of contragredience is useful throughout the theory of automorphic forms and group representations and in particular helps resolve the question of reducibility for $\rho(\mu_1, \mu_2)$.

\begin{lemma}\label{lem:4.17}

The contragredient of $\rho(\mu_1, \mu_2)$ is $\rho(\mu_1^{-1}, \mu_2^{-1})$.

\end{lemma}

This lemma results from a straightforward application of the definition of induced representation and the fact that the (one-dimensional) representations inducing $\rho(\mu_1, \mu_2)$ and $\rho(\mu_1^{-1}, \mu_2^{-1})$ are obviously contragredient to one another.

\begin{theorem}\label{thm:4.18}

The representation $\rho(\mu_1, \mu_2)$ is irreducible except when $\mu(x) = \mu_1\mu_2^{-1}(x) = |x|$ or $|x|^{-1}$. If $\mu(x) = |x|^{-1}$, then $\rho(\mu_1, \mu_2)$ contains a one-dimensional invariant subspace and the representation induced on the resulting factor space is irreducible. If $\mu(x) = |x|$, then $\rho(\mu_1, \mu_2)$ contains an irreducibly invariant subspace of codimension one. (The irreducible representations $\rho(\mu_1, \mu_2)$ are denoted by $\pi(\mu_1, \mu_2)$ and called principal series representations; if $\mu_1\mu_2^{-1}(x) = |x|$ or $|x|^{-1}$ the resulting irreducible subquotients of $\rho(\mu_1, \mu_2)$ are again denoted by $\pi(\mu_1, \mu_2)$ but are called special representations.)

\end{theorem}

\begin{proof}[Proof (Sketch)]

Roughly speaking, the idea is to analyze explicitly the ``Fourier transform'' description of a certain ``non-compact realization'' of the induced representation $\rho(\mu_1, \mu_2)$.

More precisely, recall the Bruhat decomposition for $G$. This says that

\begin{equation}\label{eq:4.11}
G = B \cup B \begin{pmatrix} 0 & -1 \\ 1 & 0 \end{pmatrix} N.\tag{4.11}
\end{equation}

(In particular, the ``big cell'' $B\begin{pmatrix} 0 & -1 \\ 1 & 0 \end{pmatrix}N$ is everywhere dense in $G$.) Now let $w = \begin{pmatrix} 0 & 1 \\ -1 & 0 \end{pmatrix}$. Then \eqref{eq:4.11} together with the transforming property \eqref{eq:4.9} implies that each $\phi \in \mathcal{H}(\mu_1, \mu_2)$ is completely determined by its restriction to $w^{-1}N$. Therefore $\mathcal{H}(\mu_1, \mu_2)$ is naturally isomorphic to the space $F(\mu_1, \mu_2)$ consisting of locally constant functions

\[
\phi(x) = \phi\!\left(w^{-1}\begin{pmatrix} 1 & x \\ 0 & 1 \end{pmatrix}\right)
\]

on $F$ with the property that $\mu(x)|x|\phi(x)$ is constant for $|x|$ large.

Now for each $\phi \in F(\mu_1, \mu_2)$ consider the ``twisted'' Fourier transform

\begin{equation}\label{eq:4.12}
\hat{\phi}(x) = \mu_2(x)|x|^{1/2}\int_F \phi(y)r(xy)\dd y.\tag{4.12}
\end{equation}

Here $r(x)$ denotes any fixed non-trivial character of the additive group of $F$. Some $p$-adic analysis shows that the integral in \eqref{eq:4.12} converges in some appropriate sense and that the resulting map

\begin{equation}\label{eq:4.13}
\phi \to \hat{\phi}\tag{4.13}
\end{equation}

is injective unless $\mu(x) = |x|^{-1}$ (in which case its kernel is the one-dimensional subspace of $\mathcal{H}(\mu_1, \mu_2)$ spanned by $\phi_0(g) = \mu_1(\det g)|\det g|^{1/2}$). On the other hand, the image space of the map \eqref{eq:4.13}, call it $\mathcal{K}(\mu_1, \mu_2)$, consists of locally constant functions on $F^{\times}$ which vanish when $|x|$ is large, and a trivial formal computation shows that

\begin{equation}\label{eq:4.14}
\hat{\phi}_*(x) = r(bx)\hat{\phi}(ax)\tag{4.14}
\end{equation}

if

\[
\phi_*(g) = \left(\rho(\mu_1, \mu_2)\left(\begin{pmatrix} a & b \\ 0 & 1 \end{pmatrix}\right)\phi\right)(g).
\]

That is, on $\mathcal{K}(\mu_1, \mu_2)$ the restriction of $\rho(\mu_1, \mu_2)$ to the subgroup $\left\{\begin{pmatrix} a & b \\ 0 & 1 \end{pmatrix}\right\}$ has a particularly simple description.

Suppose, finally, that $\mathcal{H}(\mu_1, \mu_2)$ contains some non-trivial invariant subspace $V$. The corresponding invariant subspace of $\mathcal{K}(\mu_1, \mu_2)$ will then contain some non-zero function $\hat{\phi}(x)$, hence $\hat{\phi}(x) - r(bx)\hat{\phi}(x)$ for all $b \in F$ (by \eqref{eq:4.14}), hence some non-zero function in $S(F^{\times})$, the Schwartz-Bruhat space of locally constant compactly supported functions on $F^{\times}$. (For each $b$ in $F$ the function $1 - r(bx)$ vanishes in some neighborhood of the origin.) But now $S(F^{\times})$ can easily be shown to be irreducible under the action of the operators $\hat{\phi}(x) \to r(bx)\hat{\phi}(ax)$. Therefore any non-trivial invariant subspace of $\mathcal{K}(\mu_1, \mu_2)$ contains the function $\hat{\phi}(x) - r(bx)\hat{\phi}(x)$ for every $\hat{\phi}$ in $\mathcal{K}(\mu_1, \mu_2)$, or equivalently, the invariant subspace $V$ of $\mathcal{H}(\mu_1, \mu_2)$ contains $\phi - \rho(\mu_1, \mu_2)\left(\begin{pmatrix} 1 & b \\ 0 & 1 \end{pmatrix}\right)\phi$ for all $\phi$ in $\mathcal{H}(\mu_1, \mu_2)$ (assuming $\mu(x) \neq |x|^{-1}$). So by Lemma~\ref{lem:4.17} the subspace $V^{\perp}$ orthogonal to $V$ must consist of functions $\psi(g)$ in $\mathcal{H}(\mu_1^{-1}, \mu_2^{-1})$ which are invariant under all $\rho\left(\begin{pmatrix} 1 & b \\ 0 & 1 \end{pmatrix}\right)$, $b \in F$. (Indeed $\psi \in V^{\perp}$ implies $0 = \langle \phi, \psi \rangle - \langle \rho(\mu_1, \mu_2)\left(\begin{pmatrix} 1 & b \\ 0 & 1 \end{pmatrix}\right)\phi, \psi \rangle$ or $\langle \phi, \psi \rangle = \langle \phi, \rho(\mu_1^{-1}, \mu_2^{-1})\left(\begin{pmatrix} 1 & b \\ 0 & 1 \end{pmatrix}\right)\psi \rangle$ for all $\phi$ in $\mathcal{H}(\mu_1, \mu_2)$ and $b \in F$.) Since by \eqref{eq:4.9} and \eqref{eq:4.11} such functions $\psi(g)$ do not exist (unless $\mu(x) = |x|$, in which case $\psi(g)$ must be $\mu_1^{-1}(\det g)|\det g|^{1/2}$) the proof is complete. (Actually there remains the case $\mu(x) = |x|^{-1}$, but by Lemma~\ref{lem:4.17} this situation is contragredient to the case $\mu(x) = |x|$ just considered.)

\end{proof}

\begin{remark}\label{rem:4.19}

The principal series representations $\pi(\mu_1, \mu_2)$ and $\pi(\nu_1, \nu_2)$ are equivalent if and only if $(\mu_1, \mu_2) = (\nu_1, \nu_2)$ or $(\nu_2, \nu_1)$. It is also clear from Lemma~\ref{lem:4.17} and the proof of Theorem~\ref{thm:4.18} that the special representations $\pi(\mu_1, \mu_2)$ and $\pi(\mu_2, \mu_1)$ are equivalent.

Given any fixed character $\chi(x)$, it should be noted that there are only finitely many special representations having $\chi$ as central character. (Observe that the central character of any subquotient of $\rho(\mu_1, \mu_2)$ is $\mu_1\mu_2$.)

Now the question remains: Which irreducible admissible representations of $G$ (if any) are not accounted for by the induced representations $\rho(\mu_1, \mu_2)$ and their subquotients?

\end{remark}

\begin{definition}\label{def:4.20}

Given an admissible representation $\pi$ of $G$ on $V$ let $V(\pi, N)$ denote the subspace of vectors $v$ in $V$ such that

\begin{equation}\label{eq:4.15}
\int_U \pi(n)v\dd n = 0\tag{4.15}
\end{equation}

for some open compact subgroup $U$ of $N$. Then $\pi$ is called supercuspidal if $V(\pi, N) = V$.

\end{definition}

It is well-known that the matrix coefficients

\[
(\pi(g)v, \tilde{v}) \qquad v \in V, \ \tilde{v} \in \tilde{V},
\]

of supercuspidal representations of $G$ are compactly supported (modulo the center $Z$) and that this property characterizes the supercuspidal representations among the irreducible admissible representations of $G$. Equally well known is the following fact, which is a special case of a general theorem of Jacquet and Harish Chandra.

\begin{theorem}\label{thm:4.21}

Suppose $\pi$ is an irreducible admissible representation of $G$. Then if $\pi$ is not supercuspidal it is a subrepresentation of some $\rho(\mu_1, \mu_2)$ (i.e.\ it is equivalent to some principal series or special representation). On the other hand, if $\pi$ is supercuspidal, it is not equivalent to a subquotient of any $\rho(\mu_1, \mu_2)$.

\end{theorem}

\begin{proof}[Proof (Sketch)]

Suppose $\pi$ is not supercuspidal, i.e.\ $V(\pi, N)$ is a proper subspace of $V$. Observe then that $B$ leaves $V(\pi, N)$ invariant and $N$ operates trivially on $V/V(\pi, N)$. Therefore $\pi$ induces a representation $\sigma'$ of $B/N (=A)$ in $V/V(\pi, N)$. But by the irreducibility of $\pi$ it follows that $V/V(\pi, N)$ is a $B$-module of finite type. Thus there is an invariant subspace $V'$ of $V$ which contains $V(\pi, N)$ and is such that the resulting representation $\sigma$ of $B/N$ on $W = V/V'$ is irreducible (and hence one dimensional). Since the linear map

\[
v \to f_v(g) = \underline{\pi(g)v},
\]

where $\underline{v}$ denotes the class of $v$ in $W = V/V'$, is easily seen to intertwine $(\pi, V)$ with the induced representation

\[
\Ind_B^G \sigma,
\]

and since $\sigma$ must be one-dimensional, the proof is complete.

\end{proof}

The question still remains: Do supercuspidal representations actually exist? The best way to show that they do is to construct them. As it turns out there are a number of different ways to do this. One involves the very useful notion of ``Weil representation.'' Since this construction is particularly well suited to applications to automorphic forms we present it in Section~\ref{ch:7} and again (from a different point of view) in Section~\ref{ch:10}.

\subsection{Some Properties of Irreducible Admissible Representations}

Here we collect some facts relating to the decomposition of the restriction to $K$ of an irreducible admissible representation of $G$. Of course one knows a priori that any such decomposition contains a given irreducible of $K$ at most finitely many times. But it is a more delicate question to ask exactly how many times.

\begin{definition}\label{def:4.22}

An irreducible admissible representation $\pi$ of $G$ is called class 1 or spherical if its restriction to $K$ contains the identity representation at least once.

\end{definition}

\begin{theorem}\label{thm:4.23}

An (infinite-dimensional) irreducible admissible representation $\pi$ of $G$ is class 1 if and only if $\pi = \pi(\mu_1, \mu_2)$ for some pair of unramified characters $\mu_1, \mu_2$ of $F^{\times}$ and $\pi$ is not a special representation. In this case the identity representation is contained exactly once in $\pi$. Furthermore, if $\phi_0(g)$ denotes any function in the (1-dimensional) subspace of $K$-invariant vectors in $\mathcal{H}(\mu_1, \mu_2)$, and $\tau_{\widetilde{\omega}}$ denotes the Hecke operator corresponding to convolution over $G$ with the characteristic function of the double coset $K\begin{pmatrix} \widetilde{\omega} & 0 \\ 0 & 1 \end{pmatrix}K$, then $\phi_0(g)$ is an eigenfunction of $\tau_{\widetilde{\omega}}$, and

\begin{equation}\label{eq:4.16}
\phi_0 * \tau_{\widetilde{\omega}}(g) = q^{1/2}(q^{s_1} + q^{s_2})\phi_0(g)\tag{4.16}
\end{equation}

if $\mu_i(x) = |x|^{s_i}$. (Recall that $|\widetilde{\omega}| = q^{-1}$.)

\end{theorem}

\begin{proof}

Most of the statements above are quite well-known and established in a number of places. We content ourselves with the proof of \eqref{eq:4.16}.

By definition

\[
\phi_0 * \tau_{\widetilde{\omega}}(x) = \int_G \phi(xy^{-1})\tau_{\widetilde{\omega}}(y)\dd y,
\]

so $\phi_0 * \tau_{\widetilde{\omega}}$ is again obviously right $K$-invariant. Therefore, assuming the second assertion of the theorem to be true, it follows immediately that $\phi_0$ is an eigenfunction of $\tau_{\widetilde{\omega}}$, thus in particular that $\phi_0 * \tau_{\widetilde{\omega}}(1)$ equals some scalar. But, as in Section~\ref{ch:3}, we find that

\begin{equation}\label{eq:4.17}
K\begin{pmatrix} \widetilde{\omega} & 0 \\ 0 & 1 \end{pmatrix}K = K\begin{pmatrix} \widetilde{\omega} & 0 \\ 0 & 1 \end{pmatrix} \cup \bigcup_{\eta \in O_F/P_F} K\begin{pmatrix} 1 & \eta \\ 0 & \widetilde{\omega} \end{pmatrix}.\tag{4.17}
\end{equation}

Therefore, from \eqref{eq:4.9}, \eqref{eq:4.16} and the right $K$-invariance of $\phi_0$, we have

\[
\begin{aligned}
\phi_0 * \tau_{\widetilde{\omega}}(1) &= \int_{K\begin{pmatrix} \widetilde{\omega} & 0 \\ 0 & 1 \end{pmatrix}K} \phi_0(y^{-1})\dd y\\
&= \phi_0\!\left(\begin{pmatrix} \widetilde{\omega}^{-1} & 0 \\ 0 & 1 \end{pmatrix}\right) + \sum_{\eta \in O_F/P_F} \phi_0\!\left(\begin{pmatrix} 1 & -\eta \\ 0 & \widetilde{\omega}^{-1} \end{pmatrix}\right)\\
&= q^{s_1+1/2} + q^{s_2+1/2} = q^{1/2}(q^{s_1} + q^{s_2}).
\end{aligned}
\]

\end{proof}

Even if $\pi$ is not class 1 we have the following wonderful result.

\begin{theorem}\label{thm:4.24}

Let $\pi$ denote any irreducible admissible (infinite-dimensional) representation of $G$ with central character $\psi$. Then there is a largest ideal $c(\pi)$ of $O_F$ such that the space of vectors $v$ with

\begin{equation}\label{eq:4.18}
\pi\!\left(\begin{pmatrix} a & b \\ c & d \end{pmatrix}\right)v = \psi(a)v\tag{4.18}
\end{equation}

for all

\begin{equation}\label{eq:4.19}
\begin{pmatrix} a & b \\ c & d \end{pmatrix} \in \Gamma_0'(c(\pi)) = \left\{\begin{pmatrix} a & b \\ c & d \end{pmatrix} \in K: c \equiv 0 \pmod{c(\pi)}\right\} = K^{c(\pi)}\tag{4.19}
\end{equation}

is not empty. Furthermore, this space has dimension one.

\end{theorem}

\begin{remark}\label{rem:4.25}

This theorem is proved in \cite{Casselman} and is implicit in \cite{JacquetLanglands}. Following \cite{Casselman}, we call $c(\pi)$ the conductor of $\pi$. If $\pi$ is class 1, then $c(\pi)$ is $O_F$. In general, we have the table:

\begin{equation}\label{eq:4.20}
\begin{array}{l l}
\text{Representation} & \text{Conductor}\\
\pi = \pi(\mu_1, \mu_2) & (\text{Conductor of } \mu_1)(\text{Conductor of } \mu_2)\\
\text{(principal series)} & \\
\pi = \pi(\mu_1, \mu_2) & (\text{Conductor of } \mu)^2 \text{ or } \widetilde{\omega}O_F\\
\text{(special representations)} & \\
\pi \text{ supercuspidal} & (\widetilde{\omega}^N O_F), \ N \geq 2.
\end{array}\tag{4.20}
\end{equation}

\end{remark}

Recall that the conductor of any quasi-character $\mu$ of $F^{\times}$ is the largest ideal $\widetilde{\omega}^nO_F$ such that $\mu$ is trivial on the subgroup $1 + \widetilde{\omega}^nO_F$ of $O_F^{\times}$. For the special representation $\pi(\mu_1, \mu_2)$, the conductor is $\widetilde{\omega}O_F$ if and only if $\mu$ is unramified.

As a corollary to the proof of Theorem~\ref{thm:4.24} we have that if the ideal $c(\pi)$ is replaced by the ideal $c(\pi)\widetilde{\omega}^i$ contained in it then the dimension of the space satisfying Condition \eqref{eq:4.18} for $\begin{pmatrix} a & b \\ c & d \end{pmatrix} \in \Gamma_0(c(\pi)\widetilde{\omega}^i)$ is $(i + 1)$.

Theorem~\ref{thm:4.24} is crucial to a proper understanding of the correspondence between classical cusp forms and certain representations of $\GL(2,\A)$.

\begin{definition}\label{def:4.26}

An admissible representation $\pi$ of $G$ on a complex space $V$ will be called pre-unitary (or simply unitary) if there exists on $V$ an invariant positive-definite hermitian form.

\end{definition}

If $\pi$ is pre-unitary, the operators $\pi(g)$ can be extended to unitary operators on the Hilbert space obtained by completing $V$ with respect to this form and the result is a unitary representation in the classical sense. Furthermore, this completion will be topologically irreducible if and only if $(\pi, V)$ is algebraically irreducible.

\begin{theorem}\label{thm:4.27}

The (pre-)unitary irreducible admissible representations of $G$ are:

\begin{enumerate}[(i)]

\item the supercuspidal representations with unitary central character;

\item the principal series representations $\pi(\mu_1, \mu_2)$ with both $\mu_1$ and $\mu_2$ unitary;

\item the representations $\pi(\mu_1, \mu_2)$ of the principal series for which $\mu_2(x) = \overline{\mu_1(x)}^{-1}$ and $\mu(x) = |x|^{\sigma}$, $0 < \sigma < 1$;

\item the special representations with unitary central character.

\end{enumerate}

\end{theorem}

This result is entirely analogous to the real situation (Remark~\ref{rem:4.7}). Therefore the representations in (ii) and (iii) are called continuous (respectively) complementary series representations of $G$. A proof of Theorem 4.26 may be found in \cite{JacquetLanglands} or \cite{Godement}.

\textbf{WARNING:} For $G = \GL(2,F)$ it turns out that every irreducible unitary representation of $G$ is admissible (more precisely, the corresponding representation of $G$ in the complex space of $K$-finite vectors is admissible in the sense of Definition~\ref{def:4.9}). Although it is strongly suspected that any irreducible unitary representation of any reductive $p$-adic group is admissible, at the time of this writing this remains a conjecture (in contrast to the real case).

\section{Global Theory}\label{sec:4c}

Suppose $F$ is any global field and $v$ is any place of $F$. To put together the local theory just discussed one obviously wants to be able to say that every ``nice'' representation of $\GL(2,\A_F)$ is an ``infinite tensor product'' of certain local representations $\pi_v$ of $G_v = \GL(2, F_v)$. If this tensor product is to make sense it should not seem surprising that almost all these $\pi_v$ must be class 1. But rather than continue with these motivating remarks let us simply state some facts.

For each place $v$ of $F$ let $\pi_v$ denote an irreducible unitary (hence admissible) representation of $G_v$ on $H_v$. We shall assume that the resulting collection $\{\pi_v\}$ is such that for almost every $v$ the representation $\pi_v$ is class 1. Then one can make sense out of the tensor products $\bigotimes_{v} H_v$ and $\bigotimes_{v} \pi_v$ as follows.

Let $S_0$ denote a finite set of places of $F$ containing the archimedean places as well as those where $\pi_v$ is not class 1. For each finite set $S \supset S_0$ set

\[
H_S = \bigotimes_{v \in S} H_v.
\]

For each class 1 $\pi_v$, pick once and for all a unit vector $\xi_v^0$ in the one-dimensional subspace of $K_v$-fixed vectors in $H_v$ (cf.\ Theorem~\ref{thm:4.23}). Then for each $S' \supset S$ fix an isometric imbedding of $H_S$ into $H_{S'}$, namely $\xi \to \xi \otimes \bigotimes_{v \in S' - S}\xi_v^0$. This imbedding enables us to define the (pre-)Hilbert space

\[
\lim_{\rightarrow} H_S
\]

and by completion a Hilbert space

\[
H = \bigotimes_v H_v.
\]

The resulting unitary representation

\begin{equation}\label{eq:4.21}
\pi = \bigotimes_v \pi_v\tag{4.21}
\end{equation}

of $G_\A$ in $H$ is given by

\begin{equation}\label{eq:4.22}
\pi(g)\xi = \bigotimes_v \pi_v(g_v)\xi_v\tag{4.22}
\end{equation}

if $\xi = \bigotimes_v \xi_v$ (with $\xi_v = \xi_v^0$ for almost every $v$!) and $g = (g_v)$. This representation is well-defined since $\pi_v(g_v)\xi_v = \xi_v^0$ for almost every $v$ by assumption. It is irreducible because each $\pi_v$ is. Indeed, suppose $(\cdot, \cdot)$ is a continuous hermitian bilinear form on $H$ invariant for $\pi(g)$. Then $(\cdot, \cdot)_S$, its restriction to $H_S$, is invariant for $G_S = \prod_{v \in S} G_v$, and consequently coincides with the natural inner product in $H_S$ (by Schur's Lemma). Since this holds for every finite set $S$, $(\cdot, \cdot)$ must be a multiple of the natural inner product in $H = \lim H_S$. Thus $\pi$ is irreducible.

\begin{definition}\label{def:4.28}

Suppose $\pi$ is an irreducible unitary representation of $G$ on $H$. Let $\sigma$ denote any irreducible representation of $K = \prod_v K_v$ and $H(\sigma)$ the space of vectors in $H$ which transform under $K$ according to $\sigma$. Then $\pi$ is called admissible if $H(\sigma)$ is finite dimensional for every such $\sigma$.

\end{definition}

Let us observe now that $\pi = \bigotimes_v \pi_v$ is admissible. This follows from the fact that $\sigma = \bigotimes_v \sigma_v$, with $\sigma_v$ the identity representation of $K_v$ for all but finitely many $v$, and the fact that $H(\sigma) = \bigotimes_v H_v(\sigma_v)$.

\begin{definition}\label{def:4.29}

An arbitrary irreducible unitary representation $\pi$ will be called factorizable if there are irreducible admissible unitary representations $\pi_v$ of $G_v$ such that $\pi = \bigotimes_v \pi_v$ in the sense of \eqref{eq:4.21}.

\end{definition}

\begin{theorem}\label{thm:4.30}

Every admissible irreducible unitary representation $\pi$ of $G_\A$ is factorizable, and the corresponding local components $\pi_v$ are uniquely determined by $\pi$.

\end{theorem}

\begin{remark}\label{rem:4.31}

Theorem~\ref{thm:4.30} was first proved by \cite{GelfandGraevPyatetskiiShapiro}. A more general form of it, valid for (not necessarily unitarizable) irreducible admissible representations of a certain group algebra of $G_A$, is established in \cite{JacquetLanglands}. The reason we work with unitary representations now is that (for us) a (generalized) cusp form will be an irreducible constituent of the unitary representation provided by the natural action of $G_A$ in the Hilbert space $L^2(G_F \setminus G_A)$. The usefulness of Theorem~\ref{thm:4.30} comes from the fact that any such constituent is admissible and hence can be analyzed completely in terms of its local components. (Theorem~\ref{thm:4.30} should be viewed as the analogue for $\GL(2)$ of the decomposition $\psi = \prod_v \psi_v$ valid for any unitary character of $\A^{\times} = \GL(1, \A_F)$.)

\end{remark}

\section*{Notes and References}

\addcontentsline{toc}{subsection}{Notes and References}

\textbf{A.} Our treatment of the representation theory of $\GL(2,\R)$ follows \cite{JacquetLanglands} who in turn essentially follow \cite{Bargmann} and \cite{HarishChandra2} (modulo the modifications required by the fact that $\GL(2,\R)$ is not quite the direct product of its center and $\SL(2,\R)$). In particular, the fact that every admissible representation of $\mathcal{H}(G)$ is a subquotient of some $\rho(\mu_1, \mu_2)$ follows immediately from Harish-Chandra's famous ``Subquotient Theorem'' specialized to $\SL(2,\R)$.

\textbf{B.} The representation theory of $\GL(2,F)$ over local fields was initiated by \cite{Mautner} and subsequently developed by several mathematicians (cf.\ \cite{Shalika}, \cite{Tanaka}, \cite{GelfandGraevPyatetskiiShapiro}, \cite{SallyShalika}, \cite{Silberger}, and \cite{JacquetLanglands}). Important results for $\GL(n)$ appear in \cite{Howe}.

The notion of admissibility is implicit in \cite{Mautner} and later works but was first made completely precise (and first fully exploited) in \cite{JacquetLanglands}.

In \cite{JacquetLanglands} the classification of admissible representations is accomplished through the study of Kirillov models. Our treatment (by contrast) meticulously avoids Kirillov models (except for the coincidental fact that the space $\mathcal{K}(\mu_1, \mu_2)$ which appears in the proof of Theorem~\ref{thm:4.18} provides the Kirillov model for the principal series representation $\rho(\mu_1, \mu_2)$!). Our treatment rests instead on the Subquotient Theorem~\ref{thm:4.21} (valid more generally for any reductive $p$-adic group; cf.\ \cite{HarishChandra3} and \cite{Jacquet}). The drawback of this elementary approach is that Kirillov models do play an indispensable role in the theory of automorphic forms and therefore such models have to be discussed sooner or later. (See Section~\ref{ch:6}.)
